\section{Related Work}

Drug-target affinity (DTA) prediction has been dominated by models that score a protein-drug pair directly from sequence or molecular representations. \emph{DeepDTA} established a strong sequence-based baseline by encoding protein sequences and SMILES strings \cite{weininger1988smiles} with convolutional neural networks and learning affinity from the concatenated pair representation \cite{ozturk2018deepdta}. Subsequent work improved the molecular branch while preserving the same pairwise formulation. \emph{GraphDTA} replaced the SMILES convolution with graph neural networks over molecular graphs, arguing that explicit molecular topology yields a stronger drug representation than a character sequence alone \cite{nguyen2021graphdta}. \emph{MolTrans} moved further toward learned cross-modal interaction modelling through a transformer architecture that operates on protein and compound substructures, improving benchmark performance on standard DTA splits \cite{huang2021moltrans}. Other representative models in this family include WideDTA \cite{ozturk2019widedta}, which adds protein motif and domain features, AttentionDTA \cite{zhao2023attentiondta}, which applies attention mechanisms to protein-drug pairs, and MGraphDTA \cite{yang2022mgraphdta}, which introduces multiscale graph encoders for molecular graphs. These models differ in encoder choice, but they share the assumption that the prediction problem is fundamentally pairwise: given $(p,d)$, infer affinity from that pair in isolation.

More recent work has explored alternative interaction modelling paradigms. DrugBAN \cite{bai2023drugban} introduces bilinear attention with domain adaptation to improve cross-domain DTI prediction, sharing our motivation of robustness under distribution shift. SimBoost \cite{he2017simboost} uses gradient boosting with network-based similarity features to capture relational structure between drugs and targets. Tsubaki et al.\ \cite{tsubaki2019compound} combine graph neural networks for compounds with recurrent encoders for proteins using an attention mechanism for end-to-end compound-protein interaction prediction. \emph{ConPlex} is particularly relevant to our setting because it combines a pretrained protein language model with lightweight chemical features \cite{singh2023conplex}. ConPlex constructs a contrastive co-embedding space using ProtBERT \cite{elnaggar2022prottrans} for proteins and Morgan fingerprints \cite{rogers2010morgan} for compounds, and then predicts interaction compatibility from distances in that shared latent space. This design improves efficiency and leverages large-scale protein pretraining, but it still treats affinity prediction as matching between independently encoded proteins and drugs rather than transfer from a known binding reference.

A complementary line of work uses three-dimensional structure. Structure-based docking methods can in principle model geometric complementarity more directly than sequence-only DTA models \cite{kitchen2004docking}, but they usually require reliable ligand poses or receptor conformations. DiffDock \cite{corso2023diffdock} is representative of this direction: it uses a diffusion process to generate docking poses and has shown strong performance for pose prediction, yet it depends on explicit 3D inputs for both ligand and protein and is therefore limited by structure availability and conformational uncertainty. The rapid expansion of predicted structures through AlphaFold \cite{jumper2021alphafold} and RoseTTAFold \cite{baek2021rosettafold} has made structure-aware screening more practical even for proteins lacking experimental complexes. Related protein language model systems such as ESMFold further reduce the cost of producing approximate structural models at scale \cite{lin2023esm2}. Even so, AlphaFold-based DTA pipelines typically remain structure-conditioned pairwise predictors; they address missing structure, but not the problem of transferring interaction evidence across datasets.

Protein language models have also become central to modern DTA. Large pretrained sequence encoders, including ESM-1b \cite{rives2021biological}, ProtTrans \cite{elnaggar2022prottrans} and the TAPE benchmarking framework \cite{rao2019evaluating}, have demonstrated that self-supervised pretraining on evolutionary sequences can provide richer protein features than character-level embeddings. ESM-2 in particular has emerged as a strong general-purpose protein representation model \cite{lin2023esm2}. In DTA, these embeddings are often used as frozen protein features that are combined with a drug encoder and passed to a regression or classification head \cite{li2022effective}. Our ESM-DTA baseline follows this pattern by replacing DeepDTA's learned protein embedding branch with frozen ESM-2 features while keeping the drug-side convolutional encoder. This comparison is important because it separates the value of a stronger protein representation from the value of changing the prediction formulation itself.

The evaluation practices in the DTA literature also deserve attention. Most published models are evaluated on random splits of standard benchmarks such as Davis \cite{davis2011kinase} or DTC \cite{tang2014dtc}, where significant overlap between training and test sets can inflate performance estimates. Pahikkala et al.\ \cite{pahikkala2015kronrls} showed that evaluation under cold-start (unseen protein or drug) settings is much harder, and Mayr et al.\ \cite{mayr2018deeptox} demonstrated large performance gaps across different assay collections. More recently, low-data and few-shot approaches for drug discovery have been explored \cite{altae2017low}, but these focus on learning from limited labelled examples rather than on transferring binding evidence from known interactions in a different dataset.

Taken together, prior work has substantially improved in-distribution DTA accuracy through better sequence encoders, graph encoders, transformers, contrastive co-embeddings and structure-aware docking. However, the dominant evaluation paradigm still emphasizes performance within a benchmark or on random splits of a single resource. To our knowledge, no prior DTA method explicitly frames \emph{cross-dataset transfer} as the core problem by conditioning each query on a known strong binder of the same drug. This gap motivates the present work: rather than only improving how a single protein-drug pair is encoded, we ask whether known interaction evidence can be reused as an anchor to improve generalization across datasets, protein families and difficult regimes such as intrinsically disordered proteins.
